\section{相关工作}
\paragraph{多智能体系统}
受协作式问题求解的启发，早期研究提出多智能体系统（MAS）以增强语言模型的任务求解能力~\cite{zhang2025agentorchestra, DBLP:journals/corr/abs-2505-07437,wu2024autogen, shi2025aime, gao2025single, zhu2025oagents, fang2025cognitive, zhang2024aflow,DBLP:conf/icml/Li0FXC0L25}。例如，MetaGPT~\cite{hong2023metagpt} 将智能体组织为结构化的软件开发工作流，由产品经理、架构师等专门角色通过预定义通信协议协作。OWL~\cite{hu2025owl} 采用规划器与工作器工作流，将领域无关的规划和领域相关的执行模块化，以改善迁移与泛化。尽管这类方法有效，大多数多智能体系统通常依赖固定工作流，因此较为僵化。AutoAgents~\cite{chen2023autoagents} 虽然提出为不同任务构建不同的多智能体系统，但完成任务时仍然依赖固定工作流。这促使研究逐渐转向\emph{子智能体即工具}范式，下一部分将介绍相关工作~\cite{gao2025evolvesurvey, gao2025single}。\our{} 也遵循这一范式，并进一步强调以编排为中心、动态创建子智能体，而不依赖特定的人工设计工作流。

\paragraph{子智能体即工具}
这类方法由主智能体模型以类似工具调用的方式调用子智能体来解决问题~\cite{li2025chain, su2025toolorchestra, grand2025self,DBLP:journals/tkde/LiuSLMJZFLTL25}。例如，THREAD~\cite{schroeder2025thread} 可以递归创建子智能体来处理分解后的子问题。类似地，Context-Folding~\cite{sun2025scaling} 针对子任务创建分支，再将中间步骤压缩为简洁摘要并折叠回主上下文，以此管理上下文。然而，这些方法没有把子智能体视为完全专门化的智能体，因而未能充分利用它们。Claude Code~\cite{anthropic_claude_code_subagents} 等实用系统支持在隔离上下文窗口中运行、具有自定义系统提示与工具权限的子智能体，但这些子智能体通常仍被配置为固定专家，需要人工设计。\our{} 将每个子智能体视为动态单元，并提出一种以编排为中心的智能体系统，主动按需动态创建这类子智能体，从而解决上述局限。
